\chapter{Redes Neuronales y FHE}

En el contexto de este trabajo, nos centramos en la inferencia homomórfica, donde un cliente envía sus datos cifrados a un servidor que posee el modelo (pesos sin cifrar). El servidor realiza el cómputo y devuelve el resultado cifrado. Para que este proceso es necesario redefinir las capas lineales y activaciones para que sean compatibles con la estructura algebraica del esquema CKKS.

\section{Arquitectura}

Una red neuronal puede definirse como un grafo acíclico dirigido compuesto por una secuencia de capas, donde cada capa realiza una transformación específica sobre los datos de entrada. El objetivo fundamental de esta arquitectura es aproximar una función mediante la composición de funciones más simples.

La arquitectura de una red neuronal consta de una capa de entrada, una o más capas ocultas y una capa de salida. La transición de datos entre una capa $l$ y la siguiente capa $l+1$ se define mediante una combinación de una transformación lineal seguida de una función de activación no lineal. Si representamos las activaciones de la capa $l$ como un vector $a^{l}$, el cálculo de las activaciones en la capa siguiente $a^{l+1}$ se expresa mediante la siguiente ecuación:$ a^{l+1} = \sigma(W^{l}a^{l} + b^{l}) $
Donde:

\begin{itemize}
    \item $W^{l}$ es la \textbf{matriz de pesos} de la capa $l$.
    \item $b^{l}$ es el \textbf{sesgo}, que permite desplazar la función de activación para mejorar la capacidad de ajuste del modelo.
    \item $\sigma$ es la \textbf{función de activación}, una función no lineal (ReLU, Sigmoide, etc).
\end{itemize}

\subsection{Propagación hacia Adelante}

El proceso de propagación hacia adelante es el flujo de datos desde la entrada hasta la salida para generar una predicción. En este proceso, la salida de una capa se convierte en la entrada de la siguiente. Para una red de $L$ capas, la inferencia completa es el resultado de la composición de funciones:

$$ \hat{y} = f_L(f_{L-1}(\dots f_1(x) \dots)) $

\section{El Problema de las Funciones de Activación en FHE}

En el diseño de redes neuronales convencionales, las funciones de activación no lineales son esenciales para permitir que el modelo aprenda fronteras de decisión complejas. Sin embargo, en el dominio del cifrado homomórfico, estas funciones representan el principal obstáculo para la ejecución de la inferencia.

\subsection{Limitación Algebraica de los Esquemas de Cifrado}

Los esquemas de cifrado homomórfico basados en el problema RLWE, como CKKS, se basan en estructuras de anillos matemáticos. Esto implica que las únicas operaciones aritméticas soportadas de forma nativa sobre los textos cifrados son la adición y la multiplicación.Desde una perspectiva de circuitos, esto significa que solo es posible evaluar funciones que puedan expresarse como polinomios. Cualquier operación que no sea una suma o un producto (como divisiones, comparaciones lógicas o funciones trascendentes) no puede ejecutarse directamente sobre los datos cifrados.

\subsection{Incompatibilidad de Funciones Clásicas: ReLU y Sigmoide}

Las funciones de activación más utilizadas en la actualidad presentan incompatibilidades críticas con el esquema CKKS:

\begin{itemize}
    \item \textbf{ReLU (Rectified Linear Unit):} Definida como $f(x) = \max(0, x)$. Esta función es "no suave" debido a la discontinuidad en su derivada en $x=0$. Evaluar un ReLU requiere una comparación lógica (un \textit{if-else}), operación que no existe en el álgebra de anillos de CKKS. A diferencia de esquemas como TFHE, que permiten evaluaciones de puertas lógicas, en CKKS no es eficiente realizar comparaciones de magnitud sobre el ruido del cifrado.
    \item \textbf{Sigmoide y Softmax:} Estas funciones dependen de operaciones de exponenciación ($e^x$) y división. Al no ser funciones polinomiales de grado finito, no existe una combinación de sumas y multiplicaciones homomórficas que pueda calcularlas con exactitud.
\end{itemize}

\subsection{Solución: Aproximación Polinomial}

Para superar esta limitación, la estrategia adoptada en esta tesis y en la librería desarrollada consiste en sustituir las funciones de activación no lineales por sus aproximaciones polinomiales. Este proceso transforma una función no compatible en un polinomio de grado $d$ de la forma:$ \sigma(x) \approx P_d(x) = \sum_{i=0}^{d} c_i x^i $
Existen dos métodos principales para derivar estos polinomios:

\begin{enumerate}
    \item \textbf{Polinomios de Taylor:} Útiles para aproximaciones locales alrededor de un punto, aunque su error crece rápidamente a medida que $x$ se aleja del centro de expansión.
    \item \textbf{Polinomios de Chebyshev:} Son la solución preferida en FHE y la implementada en este trabajo. Los polinomios de Chebyshev minimizan el error máximo de aproximación (error minimax) en un intervalo determinado $[-B, B]$. Esto es crucial, ya que permite mantener una precisión uniforme sobre el rango esperado de las activaciones de la red.
\end{enumerate}

La elección del grado del polinomio $d$ representa un compromiso (trade-off) técnico en la librería: un grado mayor ofrece una mejor aproximación a la función original (mayor precisión del modelo), pero consume más niveles en la jerarquía de módulos del esquema CKKS y aumenta el tiempo de cómputo debido al mayor número de multiplicaciones homomórficas requeridas.

\section{Operaciones Lineales y Empaquetado SIMD}

En el aprendizaje automático convencional, las operaciones de álgebra lineal se ejecutan mediante bibliotecas altamente optimizadas que aprovechan la memoria caché y las unidades de punto flotante de la CPU o GPU. Sin embargo, en el dominio cifrado, el desafío no es solo realizar el cálculo, sino cómo organizar los datos dentro de la estructura de "ranuras" (slots) de los textos cifrados de CKKS para maximizar el paralelismo.

\subsection{Estrategias de Layout y Empaquetado de Datos}

Como se explicó en capítulos anteriores, un texto cifrado de CKKS puede almacenar un vector de $n$ valores. Para una red neuronal, esto significa que debemos decidir cómo mapear las matrices de pesos ($W$) y los vectores de activación ($a$) a estos vectores de slots. Existen tres estrategias principales de disposición (layout):

\begin{itemize}
    \item \textbf{Layout por Filas/Columnas:} Se almacena cada fila o columna de la matriz en un texto cifrado independiente. Aunque es intuitivo, resulta ineficiente para matrices grandes, ya que consume muchos textos cifrados y aumenta el costo de comunicación.
    \item \textbf{Layout Diagonal (Método de Halevi-Shoup):} Es la técnica más eficiente para la multiplicación matriz-vector en FHE. Consiste en aplanar las diagonales de la matriz en vectores que luego se cifran. Esta estrategia permite realizar la operación completa utilizando un número mínimo de rotaciones y multiplicaciones, aprovechando al máximo la estructura SIMD.
\end{itemize}

\subsection{Multiplicación Matriz-Vector Homomórfica}

La operación fundamental de una capa densa, $y = Wx$, requiere una logística compleja en el dominio cifrado. Dado que el vector $x$ está cifrado y la matriz $W$ suele ser conocida por el servidor (texto claro), la multiplicación se realiza mediante un ciclo de Multiplicación, Rotación y Acumulación.Para multiplicar una matriz por un vector cifrado de manera eficiente, el proceso sigue estos pasos técnicos:

\begin{enumerate}
    \item \textbf{Multiplicación Elemento a Elemento:} Se multiplica el texto cifrado del vector de entrada por los vectores que representan las diagonales de la matriz de pesos (codificados en texto claro).
    \item \textbf{Rotación Cifrada (Cyclic Shift):} Debido a que los datos están empaquetados en slots, para sumar los componentes correctos es necesario rotar cíclicamente los elementos dentro del texto cifrado. Esta es una operación criptográfica específica que permite mover el valor del slot $i$ al slot $j$ sin descifrar el contenido.
    \item \textbf{Acumulación:} Se suman los resultados de las rotaciones para obtener el vector de salida cifrado.
\end{enumerate}

\subsection{El Costo de las Rotaciones}

Un aspecto crítico para el desarrollo de la librería en Python es que, en CKKS, las rotaciones son computacionalmente más costosas que las adiciones y requieren "claves de evaluación" (Galois Keys) adicionales. Por lo tanto, el algoritmo de multiplicación implementado en la librería debe optimizar el número de rotaciones necesarias. Un mal diseño del empaquetado de la matriz puede resultar en un rendimiento órdenes de magnitud más lento.Al automatizar estas estrategias de empaquetado y rotación, la librería desarrollada permite al usuario final definir capas lineales sin tener que gestionar manualmente la compleja lógica de alineación de slots y generación de claves de rotación.
